% -------------------------------------------------------------------------

\subsection{Теоретический вопрос I.1: Почему координаты в базисе единственны?}
\textbf{Формулировка вопроса:} Докажите, почему координаты любого вектора в фиксированном базисе векторного пространства
определены однозначно.

\textbf{Ответ и строгое доказательство:}
Пусть $V$ — вещественное векторное пространство, и $\mathcal{B} = (e_1, e_2, \dots, e_n)$ — базис в $V$.
По определению базиса выполнены два фундаментальных условия:
\begin{enumerate}
    \item Система векторов $(e_1, \dots, e_n)$ \textbf{порождает} все пространство: $\spanop(e_1, \dots, e_n) = V$.
    \item Система векторов $(e_1, \dots, e_n)$ \textbf{линейно независима}: равенство $\sum_{i=1}^n c_i e_i = 0$
    возможно тогда и только тогда, когда $c_1 = c_2 = \dots = c_n = 0$.
\end{enumerate}

Возьмём произвольный вектор $x \in V$.
Из первого свойства следует, что $x$ представим в виде линейной комбинации базисных векторов:
\[ x = x_1 e_1 + x_2 e_2 + \dots + x_n e_n = \sum_{i=1}^n x_i e_i, \quad x_i \in \mathbb{R}. \]
Числа $(x_1, \dots, x_n)$ называются координатами вектора $x$ в базисе $\mathcal{B}$.

\textbf{Доказательство единственности (от противного):} Предположим, что для того же самого вектора $x$ существует
второй набор координат $(y_1, y_2, \dots, y_n)$, отличный от первого:
\[ x = y_1 e_1 + y_2 e_2 + \dots + y_n e_n = \sum_{i=1}^n y_i e_i. \]
Вычтем второе векторное равенство из первого:
\[ 0 = x - x = \sum_{i=1}^n x_i e_i - \sum_{i=1}^n y_i e_i = \sum_{i=1}^n (x_i - y_i) e_i. \]
Мы получили линейную комбинацию базисных векторов $e_1, \dots, e_n$, равную нулевому вектору $0$, с коэффициентами $c_i
= (x_i - y_i)$.
Но в силу второго свойства базиса векторы $e_1, \dots, e_n$ \textbf{линейно независимы}.
Следовательно, эта линейная комбинация обязана быть тривиальной, то есть все коэффициенты равны нулю:
\[ x_i - y_i = 0 \iff x_i = y_i \quad \forall i \in \{1, 2, \dots, n\}. \]
Таким образом, $x_1 = y_1, x_2 = y_2, \dots, x_n = y_n$, что доказывает абсолютную единственность координат.

\textbf{Концептуальный вывод:} Единственность координат гарантируется свойством \emph{линейной независимости}.
Если бы базисные векторы были зависимы, существовало бы ненулевое ядро отображения координат, что привело бы к
неединственности разложения.

% -------------------------------------------------------------------------

\subsection{Теоретический вопрос I.2: Выведите формулу проекции из ортогональности}
\textbf{Формулировка вопроса:} Выведите формулу ортогональной проекции вектора на вектор и обобщите ее на случай
проекции на подпространство, опираясь исключительно на геометрическое условие ортогональности ошибки.

\textbf{Ответ и строгий вывод:}

\textbf{1.
Проекция вектора на вектор (одномерный случай):} Пусть даны ненулевой вектор $u \in V$ ($u \ne 0$) и вектор $v \in V$.
Мы ищем ортогональную проекцию вектора $v$ на прямую $\spanop(u)$, то есть вектор $p = P_u(v)$.
\begin{enumerate}
    \item Так как $p$ лежит на прямой, натянутой на $u$, он коллинеарен $u$:
    \[ p = \alpha u \quad \text{для некоторого неизвестного скаляра } \alpha \in \mathbb{R}. \]
    \item Определим вектор невязки (ошибки проекции):
    \[ e \coloneqq v - p = v - \alpha u. \]
    \item По геометрическому смыслу ортогональной проекции вектор ошибки $e$ должен быть строго \textbf{ортогонален}
    направлению проецирования $u$:
    \[ e \perp u \iff \langle e, u \rangle = 0 \iff \langle v - \alpha u, u \rangle = 0. \]
    \item Используя билинейность скалярного произведения:
    \[ \langle v, u \rangle - \alpha \langle u, u \rangle = 0 \iff \alpha \langle u, u \rangle = \langle v,
    u \rangle.\]
    \item Так как $u \ne 0$, квадрат его нормы $\langle u, u \rangle = \|u\|^2 > 0$.
    Отсюда однозначно находим скаляр $\alpha$:
    \[ \alpha = \frac{\langle v, u \rangle}{\langle u, u \rangle} = \frac{\langle v, u \rangle}{\|u\|^2}. \]
    \item Подставляя $\alpha$ в выражение для $p$, получаем каноническую \textbf{формулу ортогональной проекции}:
    \[ P_u(v) = \frac{\langle v, u \rangle}{\langle u, u \rangle} u = \frac{\langle v, u \rangle}{\|u\|^2} u. \]
\end{enumerate}

\textbf{2.
Обобщение: Проекция на подпространство (матричный вывод):} Пусть $U \le \mathbb{R}^n$ — подпространство размерности $k$,
натянутое на столбцы матрицы $A \in \mathbb{R}^{n \times k}$ полного ранга ($\rank A = k$).
Проекция вектора $y \in \mathbb{R}^n$ на $U = \Col(A)$ представляется в виде $p = A x$ для некоторого $x \in
\mathbb{R}^k$.
Вектор ошибки равен $e = y - p = y - A x$.
Условие ортогональности проекции: вектор $e$ должен быть ортогонален каждому столбцу матрицы $A$, то есть:
\[ A^\top e = 0 \iff A^\top (y - A x) = 0 \iff A^\top A x = A^\top y. \]
Поскольку $\rank A = k$, квадратная матрица Грама $A^\top A \in \mathbb{R}^{k \times k}$ строго положительно определена
и обратима.
Находим $x$:
\[ x = (A^\top A)^{-1} A^\top y. \]
Следовательно, вектор проекции равен:
\[ p = A x = A (A^\top A)^{-1} A^\top y. \]
Отсюда получаем матрицу ортогонального проектора $P = A (A^\top A)^{-1} A^\top$, для которой автоматически выполнены
свойства идемпотентности ($P^2 = P$) и симметричности ($P^\top = P$).

% -------------------------------------------------------------------------

\subsection{Теоретический вопрос I.3:Объясните иерархию metric $\to$ norm $\to$ inner product $\to$ Euclidean space}
\textbf{Формулировка вопроса:} Объясните строгую математическую иерархию геометрических структур от метрических
пространств до евклидова пространства.
Приведите условия перехода и контрпримеры.

\textbf{Ответ:}

Иерархия строится по принципу последовательного усиления аксиоматических требований:
\[ \text{Метрическое пространство} \supset \text{Нормированное векторное пространство} \supset
\text{Предгильбертово пространство} \supset \text{Евклидово пространство}.
\]

\textbf{1.
Метрическое пространство $(M, d)$:} Множество $M$ с функцией расстояния $d(x, y)$, удовлетворяющей аксиомам:
\begin{itemize}
    \item $d(x, y) \ge 0$, причём $d(x, y) = 0 \iff x = y$;
    \item $d(x, y) = d(y, x)$ (симметрия);
    \item $d(x, z) \le d(x, y) + d(y, z)$ (неравенство треугольника).
\end{itemize}
Здесь элементы не обязательно складывать и умножать на числа — есть только расстояние.

\textbf{2.
Переход «Метрика $\to$ Норма»:} Если на множестве введена структура векторного пространства $V$ над $\mathbb{R}$, можно
ввести \textbf{норму} $\|x\|$ (длину вектора):
\begin{itemize}
    \item $\|x\| \ge 0$, причём $\|x\| = 0 \iff x = 0$;
    \item $\|\alpha x\| = |\alpha| \|x\|$ (\textbf{абсолютная однородность});
    \item $\|x + y\| \le \|x\| + \|y\|$ (\textbf{субаддитивность}).
\end{itemize}
\emph{Индукция метрики:} Всякая норма порождает метрику расстояния по формуле $d(x, y) \coloneqq \|x - y\|$.
Такая метрика инвариантна к сдвигам: $d(x+z, y+z)=d(x, y)$ и однородна: $d(\alpha x,\alpha y) = |\alpha| d(x, y)$.\\
\emph{Контрпример (почему не всякая метрика есть норма):} Дискретная метрика $d(x, y) = 1$ при $x \ne y$ не
удовлетворяет свойству однородности: $d(2x, 2y) = 1 \ne 2 d(x, y)$.

\textbf{3.
Переход «Норма $\to$ Скалярное произведение»:} Скалярное произведение $\langle x, y \rangle$ — это билинейная,
симметричная и положительно определенная форма.\\ \emph{Индукция нормы:} Всякое скалярное произведение порождает
евклидову норму:
\[ \|x\| \coloneqq \sqrt{\langle x, x \rangle}. \]
Неравенство треугольника для этой нормы гарантируется неравенством Коши-Буняковского-Шварца ($|\langle x, y \rangle| \le
\|x\| \|y\|$).\\ \emph{Критерий обратного перехода (Теорема фон Неймана-Жордана):} Норма порождается скалярным
произведением тогда и только тогда, когда для любых векторов выполнено \textbf{тождество параллелограмма}:
\[ \|x + y\|^2 + \|x - y\|^2 = 2\|x\|^2 + 2\|y\|^2. \]
\emph{Контрпример (почему не всякая норма порождается скалярным произведением):} В $\mathbb{R}^2$ норма $L_1$: $\|x\|_1
= |x_1| + |x_2|$.
Возьмём $x = (1, 0)$ и $y = (0, 1)$.
Тогда $\|x\|_1 = 1$, $\|y\|_1 = 1$.
Имеем $x+y = (1, 1)$ , $\|x+y\|_1 = 2$; $x-y = (1, -1)$, $\|x-y\|_1 = 2$.
Слева: $2^2 + 2^2 = 8$.
Справа: $2(1^2) + 2(1^2) = 4$.
Так как $8 \ne 4$, тождество параллелограмма нарушено.
Следовательно, в пространстве с нормой $L_1$ понятия угла и скалярного произведения не существует.

\textbf{4.
Евклидово пространство:} Это \textbf{конечномерное} вещественное векторное пространство со скалярным произведением.
В нём одновременно существуют все четыре согласованных понятия: скалярное произведение $\langle x, y \rangle$, евклидова
норма $\|x\| = \sqrt{\langle x, x \rangle}$, евклидово расстояние $d(x, y) = \|x - y\|$ и угол $\cos \theta =
\frac{\langle x, y \rangle}{\|x\| \|y\|}$.

\newpage